\documentclass[10pt,a4paper]{amsart}
\usepackage[margin=19mm]{geometry}
\usepackage{amsmath,amssymb,mathtools}
\usepackage[scr=rsfs,cal=euler]{mathalfa}
\usepackage{xfrac}
\usepackage[unicode,hidelinks,bookmarks=false]{hyperref}
\newcommand{\ie}{i.e.\ }
\newcommand{\cf}{cf.\ }
\newcommand{\resp}{respectively\ }
\newcommand{\Subsection}{\subsection}
\providecommand{\nobreakdash}{\nobreak\hbox{-}}
\setlength{\emergencystretch}{2em}
\multlinegap0pt
\usepackage{xfrac}
\usepackage{dsfont}
\usepackage{nicefrac}
\usepackage{amssymb}
\usepackage[normalem]{ulem}
\usepackage{enumerate}
\usepackage{tikz}
\usepackage{caption}
\usepackage{booktabs}
\newtheorem{theorem}{Theorem}
\newtheorem{myenum}{Myenum}
\newtheorem{lemma}{Lemma}
\DeclareMathOperator{\GL}{GL}
\newcommand{\N}{\mathbb{N}}
\newcommand{\R}{{\mathbb{R}}}
\newcommand{\dxi}{\:\mathrm{d}\,\xi\:}
\newcommand{\dz}{\:\mathrm{d}\,z\:}
\DeclareMathOperator{\id}{id}
\newcommand{\ip}[1]{\ensuremath{\left<\text{} \, #1 \, \text{}\right>}}
\newcommand{\norm}[1]{\ensuremath{\left\lVert\text{}#1\text{}\right\rVert}}
\title{Energy Propagation in Scattering Convolution Networks Can Be Arbitrarily Slow}
\author{Hartmut F\"uhr, Max Getter}
\date{Source: arXiv:2406.05121v2 (CC BY 4.0)}
\begin{document}
\maketitle

\noindent\emph{Source and licence.} Energy Propagation in Scattering Convolution Networks Can Be Arbitrarily Slow. Version arXiv:2406.05121v2 (CC BY 4.0), distributed by arXiv under the Creative Commons Attribution 4.0 International licence, http://creativecommons.org/licenses/by/4.0/, as recorded in the arXiv licence metadata for this exact version and shown on the abstract page https://arxiv.org/abs/2406.05121v2. Full licence notice: \texttt{licenses/arxiv-2406.05121-CC-BY-4.0.txt}.\par\smallskip

\noindent\textit{Editorial scope.} Counted unit: the article's lemma add (source0.tex lines 895--899). The article's complete original proof of it is delivered with it (source0.tex lines 901--993, one \texttt{proof} environment of 93 lines). No mathematics is written, repaired or re-derived: the statement and the proof are the article's own lines, in the article's own notation, and no retained line is substituted or re-flowed.  Retained uncounted as the source-local context the proof needs: lines 345--347; lines 736--762; lines 854--858.  Nothing was omitted from any retained span, so the package carries no omission record.  No retained line contains a citation, so the wrapper carries no bibliography and no bibliography entry.  No retained line contains a figure environment, an image-inclusion command or drawing code, so no image is delivered and the package declares no asset dependency.  Editorial formatting only, in addition: an AMS \texttt{amsart} preamble that restates the article's own declarations and macros used by the retained lines, namely the environments \texttt{theorem}, \texttt{myenum}, \texttt{lemma} on separate counters, together with the standard packages those lines need.  The retained environments print in this document as Theorem 1, Myenum 1, Lemma 1 in order of appearance, whereas the article prints the same items with the numbering of its own full document.  The complete native source of the article as distributed in the arXiv e-print for this exact version is bundled unchanged as \texttt{sources/160463/source0.tex}.
  \begin{align}\label{ass: Littlewood-Paley condition}
    |\widehat{\chi}(\xi)|^2+\sum_{\psi \in \Psi} |\widehat{\psi}(\xi)|^2=1 \quad \text{a.e. } \xi \in \R^d.
  \end{align}

  \begin{theorem}\label{thm: Arbitrarily slow convergence of W_N}
    Let $F=(f_m)_{m\in \N_0} \subset L^2(\R^d)$. Suppose that for all tolerances $(\eta_k)_{k \in \N_0} \in \R_{>0}^{\N_0}$ there exist finite sets $\Psi_k\subseteq \Psi$, $k \in \N_0$, and a subsequence $(f_{m_k})_{k \in \N_0}$ of $F$ with $m_0=0$ such that the following hold:
    \begin{myenum}
      \item[(i)] The expression $C_F:=\sup_{m \in \N_0}\norm{f_m}_2^2$ is finite, i.e., $F$ is a bounded sequence.
      \item[(ii)] We have, for all $j,k \in \N_0$ with $j<k$, 
      \begin{align*}
        \left|\ip{f_{m_k},f_{m_j}}\right|\leq \eta_k.
      \end{align*}
      \item[(iii)] We have $\Psi_{k_1}\cap \Psi_{k_2}=\emptyset$ whenever $k_1\neq k_2$. 
      \item[(iv)] We have, for every $k \in \N_0$,
      \begin{align*}
        \sum_{\psi \in \Psi\setminus\Psi_k} \norm{f_{m_k}*\psi}_2^2 \leq \eta_k.
      \end{align*}  
      \item[(v)] We have, for every $k \in \N_0$ and $0\leq j < k$,
      \begin{align*}
        \sum_{\psi \in \Psi_k} \norm{f_{m_j}*\psi}_2^2 \leq \eta_k.
      \end{align*}
      \item[(vi)] There exists $\delta>0$ such that, for all $k \in \N$, $W_k(f_{m_k})\geq 4 \delta$.
    \end{myenum}
  
    If $E=(E_N)_{N \in \N}\in \R_{>0}^\N$ is a nonincreasing null-sequence, then there exists $f_E \in L^2(\R^d)$ that satisfies
    \begin{myenum}
      \item[a)] $\norm{f_E}_2^2\leq 2 C_F \cdot (1+E_1)$,
      \item[b)] $\norm{f_E-f_0}_2^2\leq 2 C_F \cdot E_1$, and
      \item[c)] for all $N\in \N$, $W_N(f_E)\geq \delta \cdot E_N$.
    \end{myenum}
  \end{theorem}

  \begin{lemma}\label{lem: ip of f and dilation of f converges to zero}
    Let $f \in L^2(\R^d)$. If $(A_m)_{m \in \N} \in \GL_d(\R)^\N$ 
    satisfies $\lim_{m \to \infty}\sigma_{\text{min}}(A_m)=\infty$, then
    \[\lim_{m \to \infty} \left|\ip{f,D_{A_m}^2f}\right|=0.\]
  \end{lemma}

  \begin{lemma}\label{lem: Matrix dilations satisfy assumptions for approx add}
    Let $(A_m)_{m \in \N} \in \GL_d(\R)^\N$ be so that $\lim_{m \to \infty}\sigma_{\text{min}}(A_m)=\infty$. Let $f \in L^2(\R^d)$, and define $f_m:=D_{A_m}^2f$, $m \in \N$. Finally, choose $f_0 \in \{0,f\}$.
    
    Then, $F=(f_m)_{m \in \N_0}$ satisfies assumptions (i)--(v) from Theorem \ref{thm: Arbitrarily slow convergence of W_N}.
  \end{lemma}

  \begin{proof}
    The statement is trivial if $f=0$. Hence, suppose $f \in L^2(\R^d)\setminus\{0\}$ in the following. Furthermore, assumptions (i)-(v) of Theorem \ref{thm: Arbitrarily slow convergence of W_N} are clearly easier to meet if we are in the case that $f_0=0$ (instead of $f_0=f$). Therefore, we only spell out the proof for $f_0=f$. For notational convenience, we set $A_0:=I_d$, so we can write $f_0=D_{A_0}^2f$.
  
    By the normalization of the dilations, $C_F:=\sup_{m\in \N_0}\norm{f_m}_2^2=\norm{f}_2^2$ is finite, i.e., assumption (i) of Theorem \ref{thm: Arbitrarily slow convergence of W_N} is satisfied. We now have to show that for all tolerances $(\eta_k)_{k \in \N_0} \in \R_{>0}^{\N_0}$ there exist finite sets $\Psi_k\subseteq \Psi$, $k \in \N_0$, and a subsequence $(f_{m_k})_{k \in \N_0}$ of $F$ with $m_0=0$ such that assumptions (ii)-(v) of Theorem \ref{thm: Arbitrarily slow convergence of W_N} hold.
    We proceed by choosing the sets $\Psi_k\subseteq \Psi$, $k \in \N_0$, and $({m_k})_{k \in \N}\in \N^\N$ inductively.
    
    Starting with $k=0$ and $m_0=0$, we only need to show that condition (iv) can be satisfied. Since 
    \[\sum_{\psi \in \Psi} \norm{f_{m_0}*\psi}_2^2\leq \norm{f}_2^2<\infty,\] there is a finite subset $\Psi_0\subset \Psi$ such that 
    \begin{align*}
      \sum_{\psi \in \Psi\setminus\Psi_0} \norm{f_{m_0}*\psi}_2^2 \leq \eta_0.
    \end{align*} 
    For the induction step, let $k \in \N$ and assume that $\Psi_0,\ldots,\Psi_{k-1}$ and $m_0,\ldots,m_{k-1}$ are such that (ii)-(v) hold. Our goal is now to obtain $\Psi_k$ and $m_k$ with the desired properties. 
    
    It is easy to meet condition (ii): For all $m \in \N_0$ and $0\leq j <k$, we have
    \begin{align*}
      \left|\ip{f_m,f_{m_j}}\right|=\left|\ip{D_{A_{m_j}^{-1}}^2D_{A_{m}}^2f,f}\right|
      =\left|\ip{D_{A_{m}A_{m_j}^{-1}}^2f,f}\right|.
    \end{align*}
    Moreover,
    \[\sigma_{\text{min}}\left(A_{m}A_{m_j}^{-1}\right)
    \geq \sigma_{\text{min}}\left(A_m\right) \cdot \sigma_{\text{min}}\left(A_{m_j}^{-1}\right)
    =\frac{\sigma_{\text{min}}\left(A_{m}\right)}{\sigma_{\text{max}}\left(A_{m_j}\right)}
    \geq \frac{\sigma_{\text{min}}\left(A_{m}\right)}{\max_{0\leq l<k}\sigma_{\text{max}}\left(A_{m_l}\right)}.\]
    Since the latter bound is uniformly in $j$ for $0\leq j<k$, and $\lim_{m \to \infty}\sigma_{\text{min}}\left(A_{m}\right)=\infty$ by assumption, Lemma \ref{lem: ip of f and dilation of f converges to zero} implies that there exists $m_k^\prime \in \N_0$ such that $\left|\ip{f_{m_k},f_{m_j}}\right|<\eta_k$ holds for all $j \in \{0,\ldots,k-1\}$ whenever we choose $m_k$ larger than $m_k^\prime$. 
  
    The harder part is to satisfy the remaining conditions (iii)-(v) simultaneously. 
    We begin with a simple observation about the impact of matrix-dilations on the frequency content of $f$.
    For $R>r>0$, let 
    \[P_{r,R}f:=\mathcal{F}^{-1}(\mathcal{F}(f)\cdot \mathds{1}_{S_{r,R}})\] 
    be the projection of $f$ on the subspace of $L^2(\R^d)$ whose Fourier transform is supported in $S_{r,R}$. We introduce $\varepsilon \in (0,\nicefrac{1}{2})$, which we will choose later in the proof (small enough). By Lebesgue's dominated convergence theorem, $\widehat{f}$ is concentrated in a spherical shell $S_{r(\varepsilon),R(\varepsilon)}$ up to $\varepsilon$-error for some $R(\varepsilon)>r(\varepsilon)>0$, meaning that
    \[\norm{P_{r(\varepsilon),R(\varepsilon)}f}_2^2
    \geq (1-\varepsilon) \cdot \norm{f}_2^2.\] 
    Now, if $A \in \GL_d(\R)$, then $\widehat{D_A^2f}$ is concentrated in $S_{\sigma_{\text{min}}(A)\cdot r(\varepsilon),\sigma_{\text{max}}(A) \cdot R(\varepsilon)}$ up to $\varepsilon$-error, since
    \begin{align*}
      \norm{P_{\sigma_{\text{min}}(A)\cdot r(\varepsilon),\sigma_{\text{max}}(A) \cdot R(\varepsilon)}D_A^2f}_2^2 
      &= \norm{\widehat{D_A^2f}\cdot \mathds{1}_{S_{\sigma_\text{min}(A)\cdot r(\varepsilon),\sigma_{\text{max}}(A)\cdot R(\varepsilon)}}}_2^2 \\
      &=\int_{\R^d}|\det(A)|^{-1} \cdot |\widehat{f}(A^{-T}\xi)|^2 \cdot \mathds{1}_{S_{\sigma_\text{min}(A)\cdot r(\varepsilon),\sigma_{\text{max}}(A)\cdot R(\varepsilon)}}(\xi) \dxi \\
      &=\int_{\R^d}|\widehat{f}(z)|^2 \cdot \mathds{1}_{S_{\sigma_\text{min}(A)\cdot r(\varepsilon),\sigma_{\text{max}}(A)\cdot R(\varepsilon)}}(A^Tz) \dz \\
      &\geq\int_{\R^d}|\widehat{f}(z)|^2 \cdot \mathds{1}_{S_{r(\varepsilon),R(\varepsilon)}}(z) \dz \\
      &=\norm{\mathcal{F}P_{r(\varepsilon),R(\varepsilon)}f}_2^2 \geq (1-\varepsilon) \cdot \norm{f}_2^2.
    \end{align*}
    Specifically, setting 
    \begin{align*}
      r_{<k}:=\min_{0\leq j<k} \sigma_{\text{min}}(A_{m_j}) \cdot r(\varepsilon) \quad \text{and} \quad R_{<k}:=\max_{0\leq j<k} \sigma_{\text{max}}(A_{m_j}) \cdot R(\varepsilon),
    \end{align*}
    the above frequency concentration estimate entails
    \begin{align*}
      \norm{(\id_{L^2(\R^d)}-P_{r_{<k},R_{<k}})f_{m_j}}_2^2 \leq \varepsilon \cdot \norm{f}_2^2 \quad \text{for } 0\leq j <k.
    \end{align*}
    By the Littlewood-Paley condition \eqref{ass: Littlewood-Paley condition} and the continuity of the Fourier transform of the filters, there is a finite subset $\widetilde{\Psi}\subseteq \Psi$ containing $\bigcup_{j=0}^{k-1}\Psi_j$, and such that
    \begin{align}\label{inequ: (5)}
      |\widehat{\chi}(\xi)|^2+\sum_{\psi \in \widetilde{\Psi}}|\widehat{\psi}(\xi)|^2 > 1-\varepsilon \quad \text{ for all } \xi \in S_{r_{<k},R_{<k}}.
    \end{align} 
    By the Riemann-Lebesgue lemma, there exists $T\geq R_{<k}$ such that
    \begin{align}\label{inequ: (4)}
      |\widehat{\chi}(\xi)|^2+\sum_{\psi \in \widetilde{\Psi}}|\widehat{\psi}(\xi)|^2 < \varepsilon \quad \text{ whenever } |\xi| \geq T.
    \end{align}
    Since $\lim_{m \to \infty}\sigma_{\text{min}}\left(A_{m}\right)=\infty$ by assumption, there is $m_k^{\prime\prime}\in \N$ with $\sigma_{\text{min}}(A_{m}) \cdot r(\varepsilon)\geq T$ for all $m \geq m_k^{\prime\prime}$.
    Define $m_k:=\max\{m_k^\prime,m_k^{\prime\prime}\}$, $r_k:=\sigma_{\text{min}}(A_{m_k}) \cdot r(\varepsilon)$, and $R_k:=\sigma_{\text{max}}(A_{m_k}) \cdot R(\varepsilon)$. Thus,
    \begin{align*}
      \norm{(\id_{L^2(\R^d)}-P_{r_{k},R_{k}})f_{m_k}}_2^2 \leq \varepsilon \cdot \norm{f}_2^2.
    \end{align*}
    Moreover, because of \eqref{inequ: (4)}, the Littlewood-Paley condition \eqref{ass: Littlewood-Paley condition}, and the continuity of the Fourier transform of the filters, there exists a finite subset $\Psi(k,\varepsilon)\subseteq \Psi \setminus \widetilde{\Psi}$ so that
    \begin{align*}
      \sum_{\psi \in \Psi(k,\varepsilon)}|\widehat{\psi}(\xi)|^2 > 1-\varepsilon \quad \text{ for all } \xi \in S_{r_k,R_k},
    \end{align*}
    which itself entails
    \begin{align}\label{inequ: (7)}
      \sum_{\psi \in \Psi \setminus \Psi(k,\varepsilon)}|\widehat{\psi}(\xi)|^2 < \varepsilon \quad \text{ for all } \xi \in S_{r_k,R_k}.
    \end{align}
    At the same time, \eqref{inequ: (5)} and the fact that $\Psi(k,\varepsilon)\subseteq \Psi \setminus \widetilde{\Psi}$, imply
    \begin{align}\label{inequ: (6)}
      \sum_{\psi \in \Psi(k,\varepsilon)}|\widehat{\psi}(\xi)|^2 < \varepsilon \quad \text{ for all } \xi \in S_{r_{<k},R_{<k}}.
    \end{align}
    We claim that $\Psi_k:=\Psi(k,\varepsilon)$ with $\varepsilon:=\frac{\eta_k}{2\norm{f}_2^2}$ does the job. 
  
    In fact, first note that conditions (i)-(iii) are immediately satisfied by the above construction. Concerning condition (iv), we have 
      \begin{align*}
        \sum_{\psi \in \Psi\setminus\Psi_k} \norm{f_{m_k}*\psi}_2^2 
        &= \sum_{\psi \in \Psi\setminus\Psi_k} \norm{(P_{r_k,R_k}f_{m_k})*\psi}_2^2 + \sum_{\psi \in \Psi\setminus\Psi_k} \norm{((\id_{L^2(\R^d)}-P_{r_k,R_k})f_{m_k})*\psi}_2^2 \\
        &\leq \int_{\R^d} |\widehat{f_{m_k}}(\xi)|^2 \cdot \sum_{\psi \in \Psi\setminus\Psi_k} |\widehat{\psi}(\xi)|^2 \cdot \mathds{1}_{S_{r_k,R_k}}(\xi)\dxi + \norm{(\id_{L^2(\R^d)}-P_{r_k,R_k})f_{m_k}}_2^2 \\
        &\leq \varepsilon \cdot \int_{\R^d} |\widehat{f_{m_k}}(\xi)|^2 \dxi + \varepsilon \cdot \norm{f}_2^2 \\
        &= \varepsilon \cdot \norm{f_{m_k}}_2^2+\varepsilon \cdot \norm{f}_2^2 = \eta_k.
      \end{align*}  
      Finally, we obtain condition (v) by proceeding analogously to the previous step. Indeed, for $0\leq j<k$, we have
      \begin{align*}
        \sum_{\psi \in \Psi_k} \norm{f_{m_j}*\psi}_2^2 
        &= \sum_{\psi \in \Psi_k} \norm{(P_{r_{<k},R_{<k}}f_{m_j})*\psi}_2^2 + \sum_{\psi \in \Psi_k} \norm{((\id_{L^2(\R^d)}-P_{r_{<k},R_{<k}})f_{m_j})*\psi}_2^2 \\
        &\leq \int_{\R^d} |\widehat{f_{m_j}}(\xi)|^2 \cdot \sum_{\psi \in \Psi_k} |\widehat{\psi}(\xi)|^2 \cdot \mathds{1}_{S_{r_{<k},R_{<k}}}(\xi)\dxi + \norm{(\id_{L^2(\R^d)}-P_{r_{<k},R_{<k}})f_{m_j}}_2^2 \\
        &\leq \varepsilon \cdot \int_{\R^d} |\widehat{f_{m_j}}(\xi)|^2 \dxi + \varepsilon \cdot \norm{f}_2^2 \\
        &= \varepsilon \cdot \norm{f_{m_j}}_2^2+\varepsilon \cdot \norm{f}_2^2 = \eta_k.
      \end{align*}
  \end{proof}

\end{document}
